\documentclass{ximera}[font=12pt]


\title{Adding Integers: The Commutative Property}   % Use xmlatex all -s to compile when SERVE refuses to work.
\author{Kelly Stady}
\license{CC: 0}         % replace with an appropriate license, or set it in xmPreamble

\begin{document}

\begin{abstract}

\end{abstract}

\maketitle
\label{xim:AddIntCommPropActivity}


\section*{The Commutative Property of Addition}

The operation of addition is \textbf{commutative}, meaning numbers can be added in any order without changing 
the result.  For example,

\begin{align*}
    4 + 3 &= 3 + 4 \\
        7 &= 7 
\end{align*}        

When adding multiple numbers with different signs, we can use this property to work more efficiently.  Instead of 
working from left to right, simply combine the positive numbers and the negative numbers separately, then add 
those results.


\begin{example}  
    
    Use the commutative property to find the sum $-6 + 3 + (-4) + 1.$

    \begin{explanation}

    Add the negative numbers $-6$ and $-4,$ then add the positive numbers $3$ and $1.$  Then add the resulting integers
    to find the answer.

    \begin{center}
        \begin{tikzpicture}[
        scale = 1.2, transform shape,      
        >=stealth, node distance=0.6cm,
        % Define Accessible Colors locally
        posColor/.style={color=accessibleOrange},
        negColor/.style={color=accessibleGlaucous}
        ]

    % --- LINE 1: THE FULL EXPRESSION ---
    % Negative 6
    \node[negColor, font=\large, inner sep=0pt] (n1minus) at (0,0) {$\mathbf{-}$};
    \node[negColor, font=\large, inner sep=0pt] (n1digit) at (0.25,0) {$\mathbf{6}$};

    \node[font=\large] (plus1) at (0.75,0) {$\mathbf{+}$};

    % Positive 3
    \node[posColor, font=\large, inner sep=0pt] (p1digit) at (1.2,0) {$\mathbf{3}$};

    \node[font=\large] (plus2) at (1.7,0) {$\mathbf{+}$};

    % Negative 4
    \node[negColor, font=\large, inner sep=0pt] (n2parenL) at (2.2,0) {$\mathbf{(}$};
    \node[negColor, font=\large, inner sep=0pt] (n2minus) at (2.4,0) {$\mathbf{-}$};
    \node[negColor, font=\large, inner sep=0pt] (n2digit) at (2.65,0) {$\mathbf{4}$};
    \node[negColor, font=\large, inner sep=0pt] (n2parenR) at (2.9,0) {$\mathbf{)}$};

    \node[font=\large] (plus3) at (3.4,0) {$\mathbf{+}$};

    % Positive 1
    \node[posColor, font=\large, inner sep=0pt] (p2digit) at (3.85,0) {$\mathbf{1}$};

    % --- RESULTS ON RIGHT (Identical thickness via \mathbf) ---
    \node[font=\large] (equals1) at (4.5,0) {$\mathbf{=}$};

    \node[negColor, font=\large, inner sep=0pt] (res1minus) at (5.0,0) {$\mathbf{-}$};
    \node[negColor, font=\large, inner sep=0pt] (res1digit) at (5.35,0) {$\mathbf{10}$};

    \node[font=\large] (plusFinal) at (5.95,0) {$\mathbf{+}$};

    \node[posColor, font=\large, inner sep=0pt] (res2digit) at (6.4,0) {$\mathbf{4}$};

    % --- TOP HORIZONTAL BAR (Negatives) ---
    \draw[negColor, thick] (0.25, 0.45) -- (2.65, 0.45);
    \draw[->, negColor, thick] (0.25, 0.45) -- (0.25, 0.2);
    \draw[->, negColor, thick] (2.65, 0.45) -- (2.65, 0.2);
    \node[negColor, font=\footnotesize\bfseries, anchor=south] (labelNeg) at (1.45, 0.45) {Add Negatives};

    % --- BOTTOM HORIZONTAL BAR (Positives) ---
    \draw[posColor, thick] (1.2, -0.45) -- (3.85, -0.45);
    \draw[->, posColor, thick] (1.2, -0.45) -- (1.2, -0.2);
    \draw[->, posColor, thick] (3.85, -0.45) -- (3.85, -0.2);
    \node[posColor, font=\footnotesize\bfseries, anchor=north] (labelPos) at (2.525, -0.45) {Add Positives};

    % --- FLOW INDICATORS ---
    \draw[->, black!75!gray, dashed, thick, shorten >=4pt] (labelNeg.north) to [out=45, in=120] (res1digit.north);  %shorten >=4pt crops the arrow back by 4pt along its trajectory, can increase or decrease
    \draw[->, black!75!gray, dashed, thick, shorten >=4pt] (labelPos.south) to [out=315, in=240, looseness=1.15] (res2digit.south); % lossness increases the sag

    % --- LINE 2: FINAL BOLD RESULT ---
    \node[font=\large, anchor=west] (final) at (equals1.west) [yshift=-1.0cm] {$\mathbf{= \, -6}$};

\end{tikzpicture}
\xmalt{Diagram showing how to use commutative property to add negative 6 plus 3 plus negative 
4 plus 1.  A blue line connects negative 6 and negative 4 with the label ''Add Negatives''.  
An orange line connects 3 and 1 with the label ''Add Positives''.  Dashed arrows show these 
groups simplifying on the right side to negative 10 plus 4.  Below this step, the final line shows 
an equals sign and the simplified result, negative 6.}
\end{center}
\end{explanation}
\end{example}

\begin{problem}
The \textbf{Commutative Property of Addition} states that changing the order of the numbers being added does not change the final sum. 

\[ a + b = b + a \]

Which of the following correctly demonstrates this property?

\begin{multipleChoice}
    \choice{$-12 + 7 = 12 + (-7)$}
    \choice{$-12 + 7 = 7 + (-12)$}
    \choice{$-12 + 7 = -12 - 7$}
    \choice[correct]{$-12 + 7 = 7 + (-12)$}
\end{multipleChoice}

\textcolor{accessibleOrange}{\textbf{Find the sum:}} $\answer{-5}$
\end{problem}

\begin{problem}
Use the \textbf{Commutative Property of Addition} to find the missing value that makes the addition statement true.

\[ 450 + (-150) = \answer{-150} + 450 \]

\textcolor{accessibleOrange}{\textbf{Find the sum:}} $\answer{300}$
\end{problem}



\begin{problem}
Use the commutative property to find each sum.

\[ -5 + 2 + (-4) + 8  = \answer{1}  \]

\[ 3 + (-1) + 6 + (-15) + 1   = \answer{-6}  \]

\[ 7 + (-12) + 3 = \answer{-2} \] 
\end{problem}

\begin{problem}
A local grassroots organization hosts a community fundraiser to support a local food bank. 
In the morning, they receive a corporate donation of $\$400$ ($+400$). 
In the afternoon, they spend $\$150$ on purchasing fresh produce for families in need ($-150$). 

According to the \textbf{Commutative Property of Addition}, changing the order of these transactions will not alter the organization's final balance.
Which of the following equations correctly demonstrates this property by showing the afternoon transaction counted first?

\begin{multipleChoice}
    \choice[correct]{$-150 + 400 = 400 + (-150)$}
    \choice{$-150 + 400 = -400 + 150$}
    \choice{$150 + 400 = 400 + 150$}
    \choice{$150 - 400 = 400 - 150$}
\end{multipleChoice}

\textcolor{accessibleOrange}{\textbf{Find the organization's final balance:}} \$ $\answer{250}$ 
\end{problem}




\section*{Great News!}

The rules for adding and subtracting integers apply to all \textbf{real numbers}.  The integers, along with the fractions and 
decimals between them, make up the set of real numbers.  

Once you have mastered these addition rules, you will be ready for subtraction.  As you will see in the next lesson, 
\textbf{if you can add integers, you can subtract them too!}


\end{document}